% Vietnamese translation for OI Public Library
% Source-File: IOI中国国家候选队论文/2007/day1/2.王晓珂《解析一类组合游戏》.ppt
\documentclass[11pt]{article}
\usepackage{oipl-vn}

\newcommand{\SG}{\operatorname{SG}}
\newcommand{\mex}{\operatorname{mex}}
\newcommand{\xor}{\mathbin{\operatorname{xor}}}

\title{Bản trình chiếu: Phân tích một lớp trò chơi tổ hợp}
\author{Wang Xiaoke}
\date{2007}

\begin{document}
\maketitle

\origtitle{Slide companion for a class of combinatorial games}
\sourcepath{IOI China candidate team papers / 2007 / day 1 / Wang Xiaoke.ppt}

\section{Phạm vi bản dịch}

Tệp PowerPoint là bản trình bày ngắn của bài viết về trò chơi tổ hợp đã được
dịch từ PDF. Bản ghi chú này giữ lại các khái niệm chính và thứ tự triển khai
thường thấy trong slide.

\section{Mô hình trò chơi}

Bài trình bày xét các trò chơi hai người, thông tin đầy đủ, không hòa, người
chơi đi luân phiên và trò chơi kết thúc sau hữu hạn bước. Một trạng thái được
gọi là thắng nếu người đang đi có thể bảo đảm thắng, và là thua nếu mọi nước
đi hợp lệ đều dẫn tới trạng thái thắng cho đối thủ.

Mô hình đồ thị có hướng giúp chuẩn hóa bài toán. Mỗi trạng thái là một đỉnh,
mỗi nước đi là một cung. Vì trò chơi kết thúc sau hữu hạn bước, đồ thị trạng
thái không có chu trình theo hướng chơi.

\section{Hàm Sprague--Grundy}

Với một trạng thái \(x\), đặt
\[
  \SG(x)=\mex\{\SG(y)\mid y \text{ là trạng thái đi tới được từ } x\}.
\]
Khi đó \(x\) là trạng thái thua khi và chỉ khi \(\SG(x)=0\). Công thức này cho
ta cách tính quy hoạch động trên đồ thị trạng thái hoặc trên cấu trúc được rút
gọn từ bài toán.

\section{Tổng của trò chơi}

Trong tổng của nhiều trò chơi con, mỗi lượt chỉ được chọn một trò chơi con để
đi. Nếu các trò chơi con có giá trị SG lần lượt là \(g_1,g_2,\ldots,g_k\), giá
trị của tổng là
\[
  g_1 \xor g_2 \xor \cdots \xor g_k.
\]
Vì vậy trạng thái tổng là thua đúng khi phép xor trên bằng \(0\). Đây là lý do
Nim và các biến thể phân rã được có thể giải bằng cùng một khuôn mẫu.

\section{Chiến lược}

Slide hướng người học qua ba thao tác tư duy: phân rã trò chơi thành các trò
chơi con độc lập, dùng quy nạp để tìm quy luật của giá trị SG, và biến đổi
trò chơi về một mô hình quen thuộc hơn. Khi không thể liệt kê toàn bộ trạng
thái, các quy luật này thường quan trọng hơn bản thân công thức truy hồi.

\end{document}
